\section{Core data model: the exact C++ types}\label{sec:datamodel}

The placement design lives or dies on the precision of a handful of value types and virtual
functions. This section fixes them exactly: the stored \emph{intent}, the derived
\emph{landmarks}, the geometric profile that powers diagnostics, the uniform way center of mass is
read in the one shared frame, and the single in-slot storage change that carries it all. The guiding
principle established in \cref{sec:philosophy} is enforced here at the type level: the only thing a
part stores about its children is physical intent; every absolute coordinate, every center of mass,
and every overlap verdict is \emph{derived} by the resolver of \cref{sec:resolver} from these
types.

A note on the starting point, because it governs how the rest of this section reads. The working
tree is entirely pre-refactor. The new types named below --- \code{StationLink}, \code{SeatKind},
\code{Station}, \code{Pose}, and the resolver --- do not yet exist, and the file
\srcf{model/parts/Placement.h} that will house the free types is new. The virtuals
\code{stationAt}, \code{radiusOuterAt}, \code{radiusInnerAt}, \code{axialLength}, and
\code{isSolid} are \emph{not} on \code{Part} today; the only geometry-adjacent virtual
the base class currently exposes is \src{model/parts/Part.h}{176} (\cppinline|getReferenceArea()|,
default \cppinline|0.0|). The discussion below therefore describes the target types and, where it
matters for correctness, flags precisely what must be added versus what already exists.

\classfig{class-placement-types}{New \texttt{Placement.h} value types and the resolver.}{fig:placement-types}{1.0}

\Cref{fig:placement-types} shows the new value types of \srcf{model/parts/Placement.h} and how the
free-function resolver consumes them; the companion \cref{fig:part-new} (landscape, below) shows the
target \code{Part} after the geometry virtuals are added.

\subsection{The station-pair link}\label{sec:datamodel-link}

A parent--child relationship in a rigid axisymmetric stack reduces to remarkably little
information: a fractional axial station on the parent, a fractional axial station on the child, an
axial gap, and a seat kind that selects the radial compatibility check and fixes the sign of the
gap. That is the whole of \code{StationLink}. Stations are stored as fractions of \emph{live}
length and resolved against the part's length on every read, so editing a part's length moves the
seam automatically rather than silently invalidating a cached number --- the staleness class this
design exists to eliminate.

\begin{listing}[htbp]
\begin{minted}{cpp}
namespace model::part {

// How a child seats against its parent. Governs (a) the radial compatibility CHECK
// and (b) the SIGN of the gap (insertion direction).
enum class SeatKind : std::uint8_t {
   Abut,        // rim-to-rim, same radius; gap is a forward standoff (signed +z)
   NestInBore,  // child OD seats inside parent ID; gap is insertion DEPTH
                //   (child moves AFT, -z, into the bore)
   OnSurface    // child seats radially on parent's outer wall at the station;
                //   gap is an axial standoff
};

} // namespace model::part
\end{minted}
\caption{The \code{SeatKind} enum: three radial relationships, governing the radial check and the
gap sign.}\label{lst:seatkind}
\end{listing}

\begin{listing}[htbp]
\begin{minted}{cpp}
// Physical intent for one parent->child attachment. Stations are FRACTIONS of live
// length, resolved against getLength() every read, so editing a length moves the seam.
struct StationLink {
   double   parentStation01{0.0};   // 0 = aft plane, 1 = fore plane, on the PARENT
   double   childStation01 {1.0};   // 0 = aft plane, 1 = fore plane, on the CHILD (default: fore)
   double   gap{0.0};               // metres; meaning per SeatKind (standoff or depth)
   SeatKind seat{SeatKind::Abut};
   Quaternion childRot{Quaternion::Identity()}; // 6-DOF seam; identity today
};
\end{minted}
\caption{The \code{StationLink} struct: two fractional stations, an axial gap, a \code{SeatKind},
and a (currently identity) seam rotation.}\label{lst:stationlink}
\end{listing}

The two members \code{parentStation01} and \code{childStation01} are dimensionless fractions in
$[0,1]$, with $0$ the aft plane and $1$ the fore plane in the \zfwd{} convention. The \code{gap} is
in meters and its meaning is selected by \code{seat}. The \code{childRot} member is the per-seam
orientation reserved for 6-DOF (\cref{sec:sixdof}); it is \code{Quaternion::Identity()} today, so
the present behavior is bit-identical to a pure-translation tree, and it costs no schema reshape to
activate later.

\Cref{fig:gap-seatkinds} makes the \code{gap} convention concrete: for each \code{SeatKind} it shows
what a zero, positive, and negative \code{gap} means physically. The resolver of \cref{sec:resolver}
places a child by \cppinline|childOriginZ = p.z + signedGap - c.z|, with \code{signedGap} equal to
\cppinline|+gap| for \seat{Abut} and \seat{OnSurface} and \cppinline|-gap| for \seat{NestInBore}. The
sign of \code{gap} is thus a displacement of the child along \zfwd, and its admissible range is part
of the seat's contract. A positive \seat{Abut} or \seat{OnSurface} gap stands the child off forward
of its flush seat, and \seat{OnSurface} alone accepts either sign --- a fore or aft standoff along
the wall. \seat{NestInBore}'s gap is a non-negative insertion \emph{depth} that drives the child aft
into the bore; a negative depth would withdraw it back out of the mouth, and a negative \seat{Abut}
gap would symmetrically drive the child into the part it mates --- a self-intersecting placement
rather than a standoff. Those degenerate signs are out-of-contract inputs the authoring helpers
(\cref{sec:authoring}) are built not to produce, distinct from the radial and envelope violations the
diagnostics of \cref{sec:diagnostics} detect between legitimately-seated parts.

\tikzfig{gap-seatkinds}{Meaning of \code{StationLink::gap} for each \code{SeatKind}, with forward
(\zfwd) to the right; gray is the parent, colour the child. \textbf{\seat{Abut}} (top): the gap is a
signed forward standoff --- flush at \mbox{gap $=0$}, an open seam for \mbox{gap $>0$}, and a
self-intersecting placement (the child driven into the part it mates) for \mbox{gap $<0$}; the child
is drawn seated forward of its parent so a positive gap reads as separation. \textbf{\seat{NestInBore}} (middle): the gap is a
non-negative insertion depth driving the child aft into the parent bore (\code{signedGap}
$=-\,$gap) --- depth $0$ rests the child at the mouth, a positive depth nests it, and a negative
value withdraws it from the bore (not seated). \textbf{\seat{OnSurface}} (bottom): the gap is an
axial standoff along the parent wall, valid in either direction (aft for \mbox{gap $<0$}, forward for
\mbox{gap $>0$}). In every case the resolver computes \mbox{\code{childOriginZ} $=$ \code{p.z} $+$
\code{signedGap} $-$ \code{c.z}}.}{fig:gap-seatkinds}

\begin{rejected}[A rich named-datum enum]
An alternative design enumerated named attachment data --- \code{FORE}, \code{MID}, \code{AFT}, and
so on --- with a parametric escape hatch for everything in between. It was rejected as paying both
costs at once: it is \emph{too rigid} (a rail button partway up a tube has no named datum and forces
the escape hatch immediately) and \emph{too rich} (the named edges \code{FORE}/\code{MID}/\code{AFT}
are nothing more than the fractions $1$, $0.5$, and $0$). Two fractional stations in $[0,1]$ subsume
the named edges, the arbitrary mid-body station, and the escape hatch in a single uniform
representation, with nothing left to extend.
\end{rejected}

\code{SeatKind}, by contrast, stays deliberately closed. A rigid axisymmetric stack genuinely
expresses only three radial relationships: rim-to-rim of equal radius (\seat{Abut}), one part's
outer diameter seated inside another's inner diameter (\seat{NestInBore}), and a part seated
radially on another's outer wall (\seat{OnSurface}). The enum's two jobs are (a) to choose the
radial compatibility check applied at the seam --- equal radii for \seat{Abut}, offender outer
radius against host capacity for \seat{NestInBore}, wall match for \seat{OnSurface} --- and (b) to
fix the sign of the gap. In the \zfwd{} frame, \seat{Abut} and \seat{OnSurface} treat the gap as a
forward standoff ($+z$), whereas \seat{NestInBore} treats it as a non-negative insertion depth that
moves the child \emph{aft} ($-z$) into the bore. There is no fourth radial relationship to add for
the parts QtRocket models, so leaving room to extend would only invite half-defined cases.

\subsection{The resolved landmark}\label{sec:datamodel-station}

Given a fractional station, a part must report a concrete geometric landmark in its own
\zfwd{} local frame: the axial coordinate and the outer and inner radii at that station. That is the
\code{Station} struct, the geometry-derived value that is computed on demand and \emph{never}
stored.

\begin{listing}[htbp]
\begin{minted}{cpp}
struct Station {
   double z{0.0};      // axial station, +z-forward (m) = (station01 - 1) * getLength()
                       //   station01: 0 = aft (z = -length), 1 = fore (z = 0, origin)
   double rOuter{0.0}; // outer radius at z (m)
   double rInner{0.0}; // inner radius at z (m); 0 for solids
};

// On Part. Base implementation uses getLength() + the radius virtuals below, so
// symmetric parts need NO override. Clamps station01 to [0,1] (see degenerate guards).
virtual Station stationAt(double station01) const;
\end{minted}
\caption{The derived \code{Station} landmark and the base \code{stationAt} virtual.}\label{lst:station}
\end{listing}

The base implementation of \code{stationAt} is written once on \code{Part} in terms of the part's
length and the radius-profile virtuals of \cref{sec:datamodel-profile}: it maps the fraction to an
axial coordinate, \cppinline|z = (station01 - 1) * getLength()| (so $z \in [-L, 0]$, with station 1 at the
fore origin and station 0 at the aft plane), and queries the outer and inner radius at
that coordinate. Because the entire computation is expressed through those virtuals, every
symmetric part inherits a correct \code{stationAt} with no override at all. A part overrides it only
if its station-to-radius mapping is not captured by the default profile.

This is where the design must add machinery the base class does not yet have.

\begin{designrule}[\code{getLength()} is promoted to a base virtual]
There is currently \emph{no} \cppinline|getLength()| on \code{Part}. Each concrete part carries its
own non-virtual accessor --- \code{BodyTube} at \src{model/parts/BodyTube.h}{51} and
\code{ConicalNoseCone} at \src{model/parts/ConicalNoseCone.h}{53} --- and \code{HollowSphere}
(\srcf{model/parts/HollowSphere.h}) has none at all, because a sphere has no length member. The
design \emph{promotes} \cppinline|getLength()| to a base virtual so the shared \code{stationAt} can
call it polymorphically. \code{HollowSphere} must then supply a definition (its axial extent is the
diameter, $2\,r_\text{outer}$), and the existing per-part accessors become overrides. This is added
work, not a pre-existing facility.
\end{designrule}

\subsection{The extent profile and the solid-host rule}\label{sec:datamodel-profile}

Diagnostics need more than a single landmark; they need the part's radial profile along its whole
axial span. Three closed-form virtuals supply it, each defaulting to a geometrically inert part on
the base class so a part with no real cross-section needs no override.

\begin{listing}[htbp]
\begin{minted}{cpp}
// On Part. Closed-form per part type. Base returns 0 (geometrically inert).
virtual double radiusOuterAt(double zLocal) const { return 0.0; } // cone tapers; tube const
virtual double radiusInnerAt(double zLocal) const { return 0.0; } // BodyTube bore; 0 solids
virtual double axialLength()  const { return getLength(); }       // span is [-axialLength(), 0]
\end{minted}
\caption{The extent-profile virtuals: outer radius, inner radius, and axial span, in the
\zfwd{} local frame.}\label{lst:profile}
\end{listing}

The reason these are distinct from \code{stationAt} is that overlap detection samples the profile at
many stations across a span (\cref{sec:diagnostics}), so the per-$z$ radius functions must be
callable independently. \code{radiusOuterAt} returns the tapering skin for a cone, the constant
outer wall for a tube, and the bulging silhouette for a sphere; \code{radiusInnerAt} returns the
bore radius for a hollow part and zero for a solid. \code{axialLength} defaults to the promoted
\cppinline|getLength()|, which is correct for the longitudinal parts and is the natural override
point for \code{HollowSphere} (whose span is its diameter).

The subtle correctness problem is deciding, for a given host part and station, the radial budget an
intruding part is allowed to occupy. A naive ``offender outer radius versus host \emph{inner}
radius'' test is wrong for a solid cone: a solid cone has no bore, so \code{radiusInnerAt} returns
zero everywhere, and the naive test would flag \emph{every} intruder as poking through. The fix is a
single predicate that selects the right boundary by whether the host is bored or solid.

\begin{designrule}[The solid-host rule]
A host \emph{occupies} radius $[0,\ \mathrm{hostCapacity}(z)]$, where
\[
\mathrm{hostCapacity}(z) \;=\;
\begin{cases}
\code{radiusInnerAt}(z) & \text{for a bored part (the bore),}\\[2pt]
\code{radiusOuterAt}(z) & \text{for a solid part (the outer skin).}
\end{cases}
\]
An offender at station $z$ with outer radius $r_\text{off}$ \emph{fits} if and only if
$r_\text{off} \le \mathrm{hostCapacity}(z) + \mathrm{tol}$.
\end{designrule}

This one predicate is correct for both cases. For a bored tube the offender must fit inside the bore
(\code{radiusInnerAt}); for a solid cone it must fit inside the outer skin (\code{radiusOuterAt})
--- which is exactly the poke-through test, the geometric question of whether the intruder breaks
the cone's outer surface. The rule is exposed as a single non-virtual helper,
\cppinline|double Part::innerCapacityAt(double zLocal) const|, so call sites in the overlap sweep
never branch on solidity themselves.

The branch is driven by one more virtual:

\begin{minted}{cpp}
virtual bool isSolid() const { return radiusInnerAt(0.0) <= 0.0; }
\end{minted}

\noindent The base default
infers solidity from the absence of a bore, which is correct for a solid cone. Parts with a wall ---
\code{BodyTube} (whose annular wall makes the bore explicit, members at
\src{model/parts/BodyTube.h}{67}) and \code{HollowSphere} (a shell, members at
\src{model/parts/HollowSphere.h}{89}) --- override \code{isSolid} to report their wall honestly, so
\code{innerCapacityAt} routes through the bore for them.

\subsection{Degenerate-geometry guards}\label{sec:datamodel-degenerate}

The profile functions must be total: a malformed or zero-extent part is a load-time diagnostic, not
undefined behavior or a divide-by-zero. Three guards make that so.

\begin{itemize}
  \item \code{stationAt} clamps \code{station01} to $[0,1]$. A link whose fraction falls outside the
        range is reported as a load-time warning rather than producing an out-of-range station.
  \item \code{radiusOuterAt} for the cone guards the taper divide: when
        $\code{length} \le 10^{-9}$ it returns \code{baseRadius} (a degenerate zero-length disc or
        ring) instead of evaluating $1 - z/\text{length}$ as $0/0$.
  \item A zero-length part (a ring or transition) has $\code{axialLength}() = 0$, so
        $\code{stationAt}(\text{any}) \to z = 0$. Its overlap span collapses to a single station,
        which the envelope sweep of \cref{sec:diagnostics} handles as a point sample. The case is
        documented behavior, not a crash.
\end{itemize}

\subsection{Center of mass in the local frame}\label{sec:datamodel-cm}

Because every part lives in the one fore-plane-origin, \zfwd{} frame (\cref{sec:philosophy-oneframe}),
center of mass is read the same way for all of them, with no per-part adapter. A part's mid-plane sits
at $z = -L/2$ (half a length aft of the fore origin), and \code{getCenterMassOffset()}
(\src{model/parts/Part.h}{96}, returning the private \code{cm} member) reports the CM relative to that
mid-plane, in the same frame. The CM's local axial station is therefore one uniform expression,
\[
\code{cmLocalZ} \;=\; -\tfrac{L}{2} \;+\; \code{getCenterMassOffset()}.z(),
\]
evaluated identically for every part type. Because \code{getCenterMassOffset()} is \emph{relative to
mid-length}, it is independent of the origin choice --- only the mid-plane term ($-L/2$) reflects the
fore origin. There is no \code{cmStationLocal} virtual and no cone special case. The values fall out
directly:

\begin{itemize}
  \item \code{BodyTube} (\src{model/parts/BodyTube.cpp}{15}) and \code{HollowSphere}
        (\src{model/parts/HollowSphere.cpp}{15}) build a centroidal tensor about a geometric-center
        CM with zero offset, so $\code{cmLocalZ} = -L/2$.
  \item \code{ConicalNoseCone}, once its reversed-frame defect is corrected
        (\cref{sec:philosophy-oneframe}), reports $\code{getCenterMassOffset()}.z() = \bar{h} - L/2$
        ($-L/4$ for a solid cone, $-L/6$ for a shell), so $\code{cmLocalZ} = \bar{h} - L$ --- i.e.\
        $-3L/4$ for a solid cone: the CM sits $3L/4$ aft of the tip, equivalently $\bar{h}=L/4$ forward
        of the wide base, where the mass actually is.
  \item \code{FinSet} (\src{model/parts/FinSet.cpp}{52}) has its axial CM at the fin-set mass centroid
        $x_c$, but its current \code{finSetCmOffset} reports $x_c$ measured from one end (the root
        leading edge), not from mid-length --- a second latent reference defect, of the same kind as
        the cone's. It is corrected the same way: report $x_c - L/2$, so
        $\code{cmLocalZ} = -L/2 + (x_c - L/2) = x_c - L$ falls out of the uniform rule with no override
        (\code{FinSetTests} updated in lock-step, \cref{sec:migration}).
\end{itemize}

\begin{keyinsight}[CM is uniform, not special-cased]
Every part's CM is derived from $-L/2 + \code{getCenterMassOffset()}.z()$ in one shared frame. The
remedy for both the cone and the fin set is the \emph{same}: report the CM relative to mid-length, in
the shared \zfwd{} frame, so the codebase carries \emph{no} reversed frame and \emph{no} non-mid
reference. No cross-frame vector addition survives anywhere in the composition path --- closing the
historical source of the simulator/visualizer divergence (\cref{sec:philosophy}) at its root rather
than papering over it with a per-part adapter.
\end{keyinsight}

\subsection{The storage swap}\label{sec:datamodel-storage}

With the link type defined, the storage change on \code{Part} is a single in-slot struct swap.
Children are stored today at \src{model/parts/Part.h}{326} as
\cppinline|std::vector<std::tuple<std::shared_ptr<Part>, Vector3>>|, the \code{Vector3} being the
child's center of mass relative to the parent's center of mass --- the CM-to-CM coupling this design
replaces. It becomes:

\begin{listing}[htbp]
\begin{minted}{cpp}
std::vector<std::pair<std::shared_ptr<Part>, StationLink>> childParts;
\end{minted}
\caption{The storage swap: the per-child \code{Vector3} offset becomes a \code{StationLink} of
physical intent, kept in a single vector.}\label{lst:storage}
\end{listing}

Keeping a \emph{single} vector is a deliberate choice.

\begin{rejected}[Parallel \code{childParts} and \code{childLinks} arrays]
Splitting the pointers and the links into two parallel vectors was rejected as a desynchronization
hazard. Every operation that reshapes the child list --- \code{clone()},
\code{removeChildById()}, and \code{addChildPart()} --- would have to keep two arrays index-aligned,
and any future edit that touches one but not the other silently corrupts the mapping. A single
vector of \cppinline|pair<ptr, StationLink>| makes the pointer and its intent inseparable by
construction; there is no index to keep aligned.
\end{rejected}

The swap is in-slot: the storage stays one vector of pairs, the same accessor pattern reads it, and
the composition routines that destructure it (\cref{sec:resolver}) change only their second binding.
Because the second element of the pair is no longer consumed as a coordinate but as intent fed to
the resolver, the \cppinline|<tuple>| dependency the current header carries solely for this member
can be dropped once the swap lands.

\subsection{What lives on \code{Part} versus in \texttt{Placement.h}}\label{sec:datamodel-where}

The design draws a clean line between behavior that must be polymorphic on a part --- the geometry
virtuals --- and the plain value types and free functions that operate over a resolved tree.
\Cref{tab:datamodel-where} states the split.

\begin{table}[htbp]
\centering
\caption{Where each new construct lives: behavior on \code{Part}; values and resolver in the new
\srcfcap{model/parts/Placement.h}.}\label{tab:datamodel-where}
\begin{tabularx}{\linewidth}{@{}XX@{}}
\toprule
\textbf{Lives on \code{Part}} & \textbf{Lives in \texttt{Placement.h} (free / POD)} \\
\midrule
\code{stationAt}, \code{radiusOuterAt}/\code{radiusInnerAt},
\code{axialLength}, \code{isSolid} (virtuals);
promoted \code{getLength()} virtual &
\code{Station}, \code{StationLink}, \code{SeatKind}, \code{Pose} \\
\addlinespace
\code{childParts} (now \code{pair<ptr, StationLink>});
non-virtual helper \code{innerCapacityAt} &
\code{resolvePlacements(...)}, \code{sweepOverlaps(...)} \\
\addlinespace
\code{addChildPart(child, StationLink)} plus a transitional \code{Vector3} overload &
\code{OverlapDiagnostic}, \code{SolveResult} \\
\bottomrule
\end{tabularx}
\end{table}

The rationale for the division is that the virtuals encode \emph{per-part geometry} and must
dispatch on dynamic type, so they belong on \code{Part}; the value types are pure data with no
behavior, and the resolver and sweep are pure functions over a whole tree, so they belong in a free
header that both the model and the visualizer can include without dragging in the part hierarchy.
The transitional \code{Vector3} overload of \code{addChildPart} (detailed in \cref{sec:authoring}
and \cref{sec:migration}) coexists with the \code{StationLink} primary during migration so the
existing call sites and \code{.qrd} fixtures keep compiling; the present tree has exactly one
overload, the legacy \cppinline|addChildPart(std::shared_ptr<Part>, Vector3)| at
\src{model/parts/Part.h}{239}.

\begin{landscape}\classfig{class-part-new}{Target \texttt{Part} after the refactor.}{fig:part-new}{1.0}\end{landscape}

\Cref{fig:part-new} shows the target \code{Part} once these virtuals, the storage swap, and the
authoring overload are in place. The composition routines that consume the resolved placement, and
the resolver itself, are the subject of \cref{sec:resolver}.
